% ============================================================
% Chapter 17: Problem 16 — Topological Problems in Algebraic Geometry
% ============================================================

\chapter{Topological Problems in Algebraic Geometry}

% ------------------------------------------------------------
\section{Historical Context}
% ------------------------------------------------------------

In the mid-nineteenth century, \textbf{Bernhard Riemann} (1826--1866) revolutionized mathematics by showing that the study of algebraic curves---the zero sets of polynomial equations in two variables---could be transformed into the study of \emph{topological surfaces}. This insight, one of the most profound in the history of mathematics, revealed a deep connection between algebra, geometry, and topology that continues to shape modern mathematics.

\subsection{Riemann Surfaces}

Consider a smooth algebraic curve $C$ defined over the complex numbers, say the curve $y^2 = x^3 - x$ in $\mathbb{C}^2$. At first glance, this is a complicated object: for each complex value of $x$, there are generally two values of $y$ (the square roots of $x^3 - x$), and as $x$ varies over $\mathbb{C}$, the curve ``branches'' at points where $x^3 - x = 0$.

Riemann's genius was to ``unfold'' this multi-valued relationship by replacing $\mathbb{C}^2$ with a new surface---now called a \textbf{Riemann surface}---on which the curve becomes a well-defined, single-valued function. Imagine two copies of the complex plane, slit along certain curves, and glued together in a clever way: the result is a two-dimensional surface on which the algebraic curve lives as a smooth manifold.

\begin{definition}[Riemann surface]
A \textbf{Riemann surface} is a one-dimensional complex manifold---that is, a surface (in the topological sense) equipped with a complex structure at each point, allowing one to perform complex analysis locally. Every smooth algebraic curve over $\mathbb{C}$ gives rise to a compact Riemann surface, and conversely, every compact Riemann surface is algebraic (this deep result is due to Riemann and later formalized by others).
\end{definition}

The Riemann surface corresponding to $y^2 = x^3 - x$ is a \textbf{torus} (a doughnut shape), not a sphere. This is a remarkable fact: the topology of the surface---its genus---carries fundamental information about the algebraic curve.

\subsection{The Genus of a Curve}

The \textbf{genus} $g$ of a compact Riemann surface is the number of ``handles'' or ``holes'' it has:
\begin{itemize}
    \item A sphere has genus $g = 0$.
    \item A torus (doughnut) has genus $g = 1$.
    \item A surface with two holes has genus $g = 2$.
\end{itemize}

For a smooth plane curve of degree $d$ in $\mathbb{P}^2$, the genus is given by the famous \textbf{genus formula}:
\[
    g = \frac{(d-1)(d-2)}{2}.
\]
Thus a line ($d=1$) has genus 0, a conic ($d=2$) has genus 0, a cubic ($d=3$) has genus 1, and a quartic ($d=4$) has genus 3. The genus is a \emph{topological invariant}: it does not change under continuous deformations of the surface.

\subsection{The Riemann--Roch Theorem}

Riemann discovered that the genus controls the space of functions on the curve. This insight was made precise by his student Gustav Roch in the \textbf{Riemann--Roch theorem}, one of the most important results in algebraic geometry.

\begin{theorem}[Riemann--Roch]
Let $C$ be a smooth algebraic curve of genus $g$, and let $D$ be a divisor (a formal sum of points on $C$). Then
\[
    \ell(D) - \ell(K - D) = \deg(D) - g + 1
\]
where $\ell(D)$ is the dimension of the space of meromorphic functions $f$ with $(f) + D \geq 0$, $K$ is the canonical divisor (the divisor of a holomorphic 1-form), and $\deg(D)$ is the degree of $D$.
\end{theorem}

For a curve, the canonical divisor $K$ has degree $2g - 2$. When $\deg(D) > 2g - 2$, the term $\ell(K - D)$ vanishes, and the theorem gives an exact formula for $\ell(D)$. The Riemann--Roch theorem shows that the \emph{topology} (the genus) of a curve completely determines many of its algebraic properties---a stunning connection between topology and algebra.

\subsection{From Curves to Surfaces: Early Topological Questions}

Riemann's work raised natural questions that went beyond curves:

\begin{itemize}
    \item \textbf{What about surfaces?} An algebraic surface (the zero set of a polynomial in three variables, or a two-dimensional variety) has a much richer topology than a curve. What topological invariants describe a surface? What constraints does the algebraic structure impose?

    \item \textbf{What about higher dimensions?} If algebraic curves give Riemann surfaces, what topology do $n$-dimensional algebraic varieties have? How do the algebraic properties of the variety constrain its topology?

    \item \textbf{How can we compute topological invariants?} Given an algebraic variety defined by equations, how can we compute its Betti numbers, its fundamental group, or its cohomology ring?
\end{itemize}

These questions---roughly speaking---constitute Hilbert's sixteenth problem. Hilbert asked for the development of a general theory connecting algebraic geometry and topology, analogous to what Riemann had achieved for curves.

% ------------------------------------------------------------
\section{The Problem}
% ------------------------------------------------------------

\begin{problem_statement}[Hilbert's Sixteenth Problem]
Develop a topology for algebraic varieties, particularly for algebraic curves and surfaces, analogous to what has been done for manifolds. Understand the topological properties of algebraic varieties.
\end{problem_statement}

Hilbert's formulation was characteristically open-ended. The problem can be broken into several sub-questions:

\begin{enumerate}
    \item \textbf{Topological invariants of varieties.} Compute the fundamental group, Betti numbers, and cohomology ring of algebraic varieties. How do these depend on the defining equations?

    \item \textbf{Constraints from algebra.} What restrictions does the algebraic structure of a variety place on its topology? For instance, not every topological manifold can be the underlying space of an algebraic variety.

    \item \textbf{Topology of real algebraic varieties.} Understand the topology of the \emph{real} points of an algebraic variety---the set of real solutions to polynomial equations. Hilbert specifically mentioned the topology of real algebraic surfaces and curves.
\end{enumerate}

The first two questions relate to \textbf{complex algebraic geometry} and connect to the rapidly developing field of algebraic topology. The third question, about real algebraic varieties, remains one of the deepest and most difficult parts of Hilbert's problem, closely related to the topology of real manifolds and the study of real algebraic sets.

% ------------------------------------------------------------
\section{Key Developments}
% ------------------------------------------------------------

The twentieth century brought extraordinary progress on Hilbert's sixteenth problem, driven by the intertwined development of algebraic topology, algebraic geometry, and differential geometry.

\subsection{Algebraic Topology: Poincar\'e and Brouwer}

The early twentieth century saw the birth of \textbf{algebraic topology}, the study of topological spaces through algebraic invariants.

\textbf{Henri Poincar\'e} (1854--1912) introduced the \textbf{fundamental group} $\pi_1(X)$ of a topological space $X$---the group of loops based at a point, up to continuous deformation. For algebraic varieties, the fundamental group encodes how ``multiply connected'' the space is:
\begin{itemize}
    \item The sphere $\mathbb{S}^2$ is simply connected: $\pi_1(\mathbb{S}^2) = 0$.
    \item The torus $\mathbb{T}^2$ has $\pi_1(\mathbb{T}^2) \cong \mathbb{Z}^2$: there are two independent loops that cannot be contracted.
    \item A Riemann surface of genus $g$ has $\pi_1(\Sigma_g) \cong \langle a_1, b_1, \ldots, a_g, b_g \mid [a_1, b_1] \cdots [a_g, b_g] = 1 \rangle$.
\end{itemize}

\textbf{L.E.J. Brouwer} (1881--1966) introduced \textbf{simplicial homology} and proved foundational results about topological invariance. The \textbf{Betti numbers} $b_k(X) = \rank H_k(X)$ of a topological space count the number of $k$-dimensional ``holes'':
\begin{itemize}
    \item $b_0(X)$ counts connected components.
    \item $b_1(X)$ counts independent loops.
    \item $b_2(X)$ counts independent voids (enclosed cavities).
\end{itemize}

For a smooth complex projective variety of complex dimension $n$ (real dimension $2n$), the Betti numbers satisfy $b_k = b_{2n-k}$ (Poincar\'e duality), and the \textbf{Euler characteristic} $\chi(X) = \sum_{k=0}^{2n} (-1)^k b_k$ is a topological invariant that encodes much information.

\subsection{Sheaf Cohomology: Leray and Serre}

The classical tools of algebraic topology (singular homology, simplicial homology) proved insufficient for the subtle problems of algebraic geometry. A new tool was needed: \textbf{sheaf cohomology}.

\textbf{Jean Leray} (1906--1998), while a prisoner of war in Germany during World War~II, developed the theory of \textbf{spectral sequences} and sheaf cohomology---one of the most remarkable mathematical achievements of the twentieth century. \textbf{Jean-Pierre Serre} (born 1926) subsequently showed that sheaf cohomology is the right tool for algebraic geometry.

\begin{definition}[Sheaf cohomology]
Let $X$ be a topological space and $\mathcal{F}$ a sheaf of abelian groups on $X$. The \textbf{sheaf cohomology groups} $H^k(X, \mathcal{F})$ are defined as the right derived functors of the global sections functor. They measure the ``failure'' of local cohomological data to patch together globally.
\end{definition}

Serre proved that for a smooth projective variety $X$ over $\mathbb{C}$, the sheaf cohomology groups $H^k(X, \mathbb{Z})$ are \textbf{finitely generated abelian groups} whose ranks are the Betti numbers $b_k$. Moreover, the \textbf{Hodge decomposition} decomposes the complex cohomology:
\[
    H^k(X, \mathbb{C}) = \bigoplus_{p+q=k} H^{p,q}(X)
\]
where $H^{p,q}(X)$ is the space of harmonic $(p,q)$-forms, or equivalently the Dolbeault cohomology group $H^{p,q}_{\bar\partial}(X)$. The dimensions $h^{p,q} = \dim H^{p,q}(X)$ are the \textbf{Hodge numbers}, and they satisfy the symmetry $h^{p,q} = h^{q,p}$ (complex conjugation) and $h^{p,q} = h^{n-p,n-q}$ (Hodge duality). This decomposition is a profound constraint on the topology of algebraic varieties: not every topological manifold has a Hodge decomposition.

\subsection{The Lefschetz Hyperplane Theorem}

One of the most powerful tools for understanding the topology of algebraic varieties is the \textbf{Lefschetz hyperplane theorem}, proved by \textbf{Solomon Lefschetz} (1884--1972) and later generalized by others.

\begin{theorem}[Lefschetz Hyperplane Theorem]
Let $X \subset \mathbb{P}^n$ be a smooth projective variety of dimension $d$, and let $H$ be a smooth hyperplane section (i.e., $H = X \cap \{L = 0\}$ for a general linear form $L$). Then the restriction maps
\[
    H_k(H, \mathbb{Z}) \to H_k(X, \mathbb{Z})
\]
are isomorphisms for $k < d - 1$ and surjective for $k = d - 1$. Dually, the restriction maps
\[
    H^k(X, \mathbb{Z}) \to H^k(H, \mathbb{Z})
\]
are isomorphisms for $k < d - 1$ and injective for $k = d - 1$.
\end{theorem}

In intuitive terms, the theorem says that a hyperplane section of a variety ``sees'' most of the topology of the ambient variety. By slicing repeatedly with hyperplanes, one can reduce questions about the topology of a high-dimensional variety to questions about lower-dimensional sections. This inductive technique is fundamental in algebraic geometry.

\subsection{The Weil Conjectures}

Perhaps the most spectacular achievement connecting topology and algebraic geometry was the proof of the \textbf{Weil conjectures}, formulated by \textbf{Andr\'e Weil} (1906--1998) in 1949 and proved by \textbf{Alexander Grothendieck} (1928--2014) and \textbf{Pierre Deligne} (born 1944) in 1974.

The Weil conjectures concern the \textbf{zeta function} of a smooth projective variety $X$ over a finite field $\mathbb{F}_q$:
\[
    Z(X, t) = \exp\left(\sum_{n=1}^{\infty} \frac{N_n}{n} t^n\right)
\]
where $N_n = \#X(\mathbb{F}_{q^n})$ is the number of $\mathbb{F}_{q^n}$-rational points of $X$.

\begin{theorem}[Weil Conjectures -- proved]
The zeta function $Z(X, t)$ of a smooth projective variety $X$ over $\mathbb{F}_q$ is a \textbf{rational function}:
\[
Z(X, t) = \frac{P_1(t)\cdots P_{2d-1}(t)}{P_0(t)\cdots P_{2d}(t)},
\]
where each $P_i(t)$ is a polynomial with integer coefficients, $P_0(t)=1-t$, $P_{2d}(t)=1-q^d t$. Moreover:
\begin{enumerate}
    \item \textbf{Rationality:} $Z(X, t)$ is a rational function.
    \item \textbf{Functional equation:} $Z(X, 1/q^d t) = \pm q^{d\chi/2} t^{\chi} Z(X, t)$, where $\chi = \chi(X)$ is the Euler characteristic.
    \item \textbf{Riemann hypothesis:} The zeros and poles of $Z(X, t)$ have absolute value $q^{-i/2}$ for appropriate $i$.
    \item \textbf{Betti number interpretation:} The degree of $P_i$ is the $i$-th Betti number $b_i$ of $X$ (computed over $\mathbb{C}$, via the Weil cohomology theory).
\end{enumerate}
\end{theorem}

The Weil conjectures are extraordinary because they connect:
\begin{itemize}
    \item \textbf{Algebra} (counting solutions over finite fields),
    \item \textbf{Topology} (Betti numbers and Euler characteristic),
    \item \textbf{Number theory} (the distribution of rational points).
\end{itemize}

Grothendieck's contribution was to reformulate the conjectures in terms of \textbf{\'{e}tale cohomology}, a new cohomology theory for algebraic varieties defined over arbitrary fields (not just $\mathbb{C}$). The key insight was that \'etale cohomology should behave like singular cohomology but work over finite fields. Grothendieck proved rationality and the functional equation. Deligne completed the proof by establishing the deepest part---the Riemann hypothesis---using a brilliant argument involving estimates for l-adic sheaves.

The Weil conjectures demonstrate that the topology of an algebraic variety (encoded in its Betti numbers and cohomology) controls its arithmetic behavior (the number of rational points). This is a profound vindication of Hilbert's intuition that topology and algebraic geometry are deeply intertwined.

\subsection{The Hodge Conjecture}

The most famous open problem arising from Hilbert's sixteenth problem is the \textbf{Hodge conjecture}, one of the seven \textbf{Clay Millennium Problems}, carrying a prize of one million dollars.

\begin{conjecture}[Hodge Conjecture]
Let $X$ be a smooth complex projective variety of dimension $n$. Then every \textbf{Hodge class}---every element of $H^{2k}(X, \mathbb{Q}) \cap H^{k,k}(X)$---is a rational linear combination of classes of algebraic subvarieties of codimension~$k$.
\end{conjecture}

In other words, certain topological cycles on algebraic varieties (those satisfying a condition from the Hodge decomposition) must be ``algebraic'': they must come from actual algebraic subvarieties. This would mean that the Hodge decomposition---a priori a topological or analytic invariant---completely detects algebraic subvarieties.

\begin{itemize}
    \item The Hodge conjecture is \textbf{trivially true} in dimension 1: every class in $H^{1,1}$ on a curve is algebraic.
    \item For surfaces (dimension 2), the Hodge conjecture is known for classes in $H^{1,1}$ (the Lefschetz $(1,1)$-theorem, proved by Lefschetz and later generalized): every integral $(1,1)$-class on a surface is the class of a divisor.
    \item For higher codimension, the Hodge conjecture remains \textbf{wide open} in general. It is known in special cases (e.g., for abelian varieties of small dimension, and for some specific classes), but no general proof is known.
\end{itemize}

The Hodge conjecture connects the \emph{analytic} structure of a variety (the Hodge decomposition) to its \emph{algebraic} structure (the existence of algebraic subvarieties). If true, it would establish a perfect correspondence between topology and algebra for complex projective varieties.

\subsection{The Topology of Real Algebraic Varieties}

Hilbert also asked about the topology of \textbf{real algebraic varieties}---the sets of real solutions to polynomial equations. This aspect of the problem is closely related to \textbf{real algebraic geometry} and \textbf{real enumerative geometry}.

Key results include:

\begin{itemize}
    \item \textbf{Harnack's theorem} (1876): A smooth real projective curve of degree $d$ has at most $g + 1$ connected components (ovals), where $g = (d-1)(d-2)/2$ is the genus. Moreover, this bound is achieved for curves of every degree.

    \item \textbf{The Smith--Thom inequality}: If $X$ is a real algebraic variety and $X(\mathbb{C})$ is its complexification, then $\chi(X(\mathbb{R})) \leq \chi(X(\mathbb{C}))$ (the sum of the Betti numbers of the real part does not exceed that of the complex part).

    \item \textbf{Viro's patchworking method}: A powerful technique for constructing real algebraic varieties with prescribed topology, connecting tropical geometry to real algebraic geometry.

    \item \textbf{Thom's conjecture} (and its generalizations): The topology of real algebraic surfaces in $\mathbb{R}^3$ is highly constrained. For instance, the maximal number of connected components of a smooth real algebraic surface of degree $d$ in $\mathbb{R}^3$ is an active area of research.
\end{itemize}

The study of real algebraic varieties remains one of the most challenging aspects of Hilbert's problem, and many fundamental questions are still open.

% ------------------------------------------------------------
\section{Current Status: Partially Solved}
% ------------------------------------------------------------

\begin{partiallybox}
Hilbert's sixteenth problem is \textbf{partially solved}. The theory of complex algebraic varieties has seen spectacular progress: the development of sheaf cohomology (Leray, Serre, Grothendieck), the proof of the Lefschetz hyperplane theorem, the Hodge decomposition, and the proof of the Weil conjectures (Grothendieck, Deligne) have given us a deep understanding of the topology of complex algebraic varieties. The Hodge conjecture---one of the Clay Millennium Problems---remains open and is considered one of the most important unsolved problems in mathematics. The topology of real algebraic varieties is partially understood (Harnack's theorem, Viro's method, the Smith--Thom inequality), but many fundamental questions remain open, including the topology of real algebraic surfaces and the distribution of connected components of real varieties.
\end{partiallybox}

% ------------------------------------------------------------
\section*{Common Questions}

\begin{description}[font=\bfseries, leftmargin=2em, style=nextline]

\item[Q: What's the difference between an algebraic variety and a manifold?]\hfill\\
A manifold is a space that locally looks like $\mathbb{R}^n$ (smooth and without singularities). An algebraic variety is the zero set of a system of polynomial equations, and it can have \emph{singularities}---points where it is not smooth, like the tip of a cone ($x^2 + y^2 = z^2$). Over the complex numbers, a smooth algebraic variety is a complex manifold, but not all complex manifolds are algebraic. The interplay between these two viewpoints---algebraic equations and smooth topology---is at the heart of Hilbert's 16th problem.

\item[Q: What are the Weil conjectures in simple terms?]\hfill\\
The Weil conjectures (proved by Deligne in 1974) relate the number of solutions to polynomial equations over finite fields to the topological structure of the corresponding complex algebraic variety. They predicted that the generating function counting solutions modulo $p$, modulo $p^2$, etc. is a rational function whose properties mirror the topology (Betti numbers, Poincaré duality) of the variety over $\mathbb{C}$. This is astonishing: counting points over finite fields reveals the shape of a space over the complex numbers.

\item[Q: What's the Hodge conjecture about?]\hfill\\
The Hodge conjecture asks whether every closed differential form of a certain type (a $(p,p)$-form) on a smooth complex projective variety can be expressed as a combination of algebraic cycles---essentially, whether topological information can be realized geometrically. It is one of the Clay Millennium Problems and remains open. In simpler terms: can every ``nice'' cohomology class be represented by an actual algebraic subvariety? Most mathematicians believe the answer is yes, but no one has proved it.

\item[Q: Why are real algebraic varieties harder than complex ones?]\hfill\\
Over $\mathbb{C}$, the fundamental theorem of algebra guarantees that polynomial equations always have solutions, and the topology is governed by powerful tools like sheaf cohomology and Hodge theory. Over $\mathbb{R}$, a polynomial equation can have no real solutions at all, components can appear and disappear as parameters change, and the topology (number of ovals, their nesting, etc.) is much more delicate. Hilbert's 16th problem specifically asks about the topology of \emph{real} algebraic curves and surfaces, which remains only partially understood.

\item[Q: What is a Lefschetz pencil?]\hfill\\
A Lefschetz pencil is a family of hyperplane sections of a projective variety, parameterized by a line in the dual projective space. It is a central tool for studying the topology of varieties: as you vary the hyperplane, the sections change smoothly except at finitely many ``vanishing cycles,'' where the section acquires a simple singularity (a node). These vanishing cycles generate the middle-dimensional homology of the variety. Lefschetz's work established a powerful connection between the algebra of equations and the topology of their solution sets.

\item[Q: How does Hilbert's 16th problem split into two parts?]\hfill\\
Part~A concerns the topology of real algebraic curves in $\mathbb{P}^2$: how many connected components (ovals) can a real plane curve of degree $d$ have, and how can they be arranged? This is mostly settled: Harnack's theorem gives the maximum number, and the possible arrangements are known for low degrees but not in general. Part~B concerns the topology of real algebraic surfaces and higher-dimensional varieties, as well as the limit cycles of polynomial vector fields in the plane---a famously difficult problem in dynamical systems that remains largely open.

\end{description}

% ------------------------------------------------------------
\section*{Exercises}
% ------------------------------------------------------------

\begin{enumerate}

    \item \textbf{Genus of plane curves.}
    \begin{enumerate}
        \item Compute the genus of a smooth plane curve of degree $d = 4$ (a quartic) using the formula $g = (d-1)(d-2)/2$.
        \item What is the genus of a smooth plane curve of degree $d = 5$ (a quintic)?
        \item A smooth plane cubic ($d = 3$) has genus 1. Its Riemann surface is a torus. Explain why this means that a cubic curve has a group structure (the sum of two points on the curve can be defined geometrically using lines through the curve).
    \end{enumerate}

    \item \textbf{Betti numbers of surfaces.}
    \begin{enumerate}
        \item The sphere $\mathbb{S}^2$ has Betti numbers $b_0 = 1$, $b_1 = 0$, $b_2 = 1$. Compute the Euler characteristic $\chi(\mathbb{S}^2) = b_0 - b_1 + b_2$.
        \item A torus $\mathbb{T}^2$ has Betti numbers $b_0 = 1$, $b_1 = 2$, $b_2 = 1$. Compute $\chi(\mathbb{T}^2)$.
        \item A Riemann surface of genus $g$ has Betti numbers $b_0 = 1$, $b_1 = 2g$, $b_2 = 1$. Show that $\chi = 2 - 2g$. Verify this for $g = 0$ (sphere) and $g = 1$ (torus).
        \item A smooth curve of degree 4 in $\mathbb{P}^2$ has genus 3. What is its Euler characteristic?
    \end{enumerate}

    \item \textbf{Fundamental groups.}
    \begin{enumerate}
        \item The fundamental group of the sphere $\mathbb{S}^2$ is trivial: $\pi_1(\mathbb{S}^2) = 0$. Explain what this means intuitively (every loop on the sphere can be contracted to a point).
        \item The fundamental group of the torus $\mathbb{T}^2$ is $\mathbb{Z} \times \mathbb{Z}$. Describe two loops on the torus that generate this group.
        \item The fundamental group of a Riemann surface of genus $g \geq 1$ is the surface group $\langle a_1, b_1, \ldots, a_g, b_g \mid [a_1, b_1] \cdots [a_g, b_g] = 1 \rangle$. For $g = 1$, this reduces to $\langle a, b \mid aba^{-1}b^{-1} = 1 \rangle \cong \mathbb{Z}^2$, consistent with the torus.
    \end{enumerate}

    \item \textbf{The Hodge diamond.}
    \begin{enumerate}
        \item The Hodge numbers $h^{p,q} = \dim H^{p,q}(X)$ of a smooth complex projective variety are arranged in a diamond-shaped array called the \textbf{Hodge diamond}. For a smooth curve of genus $g$, the Hodge diamond is:
        \[
        \begin{array}{ccc}
         & h^{0,0} & \\
        h^{1,0} & & h^{0,1} \\
         & h^{1,1} &
        \end{array}
        = 
        \begin{array}{ccc}
         & 1 & \\
        g & & g \\
         & 1 &
        \end{array}
        \]
        Verify that $h^{p,q}$ satisfies $h^{p,q} = h^{q,p}$ (symmetry) and that $\sum_{p+q=k} h^{p,q} = b_k$ (the $k$-th Betti number).
        \item For a smooth projective surface $S$, the Hodge numbers satisfy $h^{0,0} = h^{2,2} = 1$, $h^{1,0} = h^{0,1} = h^{2,1} = h^{1,2}$, $h^{2,0} = h^{0,2}$, and $h^{1,1}$ is determined by the Euler characteristic. If $S$ is a smooth quartic surface in $\mathbb{P}^3$, show that $\chi(S) = 24$ and compute the Hodge numbers.
    \end{enumerate}

    \item \textbf{Lefschetz hyperplane theorem (intuition).}
    \begin{enumerate}
        \item Let $C$ be a smooth plane curve of degree $d$ in $\mathbb{P}^2$ (so $C$ has complex dimension 1, real dimension 2). The Lefschetz hyperplane theorem says that the map $H_0(C) \to H_0(\mathbb{P}^2)$ is an isomorphism (both spaces are connected) and $H_1(C) \to H_1(\mathbb{P}^2)$ is surjective. Since $H_1(\mathbb{P}^2) = 0$, this says $H_1(C)$ surjects onto 0---which is trivially true. The interesting content is that $H^1(C)$ carries information about $C$ that is \emph{new}: it is not detected by the ambient $\mathbb{P}^2$. Explain why this is consistent with the fact that $b_1(C) = 2g > 0$ for $g \geq 1$.
        \item For a smooth surface $S$ in $\mathbb{P}^3$, the Lefschetz theorem says $H_1(S) \cong H_1(\mathbb{P}^3) = 0$ and $H_2(S)$ surjects onto $H_2(\mathbb{P}^3) \cong \mathbb{Z}$. What does this tell you about $b_1(S)$?
    \end{enumerate}

    \item[$\star$] \textbf{The Weil conjectures and point counting.}
    \begin{enumerate}
        \item Let $X = \mathbb{P}^1$ (the projective line) over $\mathbb{F}_q$. The zeta function is $Z(\mathbb{P}^1, t) = 1 / ((1-t)(1-qt))$. Compute $N_n = \#X(\mathbb{F}_{q^n})$ from this zeta function and verify that $N_n = q^n + 1$. (This counts the points of $\mathbb{P}^1(\mathbb{F}_{q^n})$: the $q^n$ affine points plus the point at infinity.)
        \item The Betti numbers of $\mathbb{P}^1 \cong \mathbb{S}^2$ are $b_0 = 1$, $b_1 = 0$, $b_2 = 1$. The Euler characteristic is $\chi = 2$. Verify that the functional equation of the zeta function holds for $\mathbb{P}^1$.
        \item (Optional) Compute the zeta function of the elliptic curve $E: y^2 = x^3 + x$ over $\mathbb{F}_3$. Count the points $N_1 = \#E(\mathbb{F}_3)$ and $N_2 = \#E(\mathbb{F}_9)$, and verify that $Z(E, t) = (1 - at + 3t^2) / ((1-t)(1-3t))$ where $a = 3 + 1 - N_1$.
    \end{enumerate}

    \item[$\star$] \textbf{Real algebraic curves (Harnack's theorem).}
    \begin{enumerate}
        \item A smooth real conic ($d = 2$) has genus 0, so it has at most $0 + 1 = 1$ connected component. Verify: the conic $x^2 + y^2 = 1$ in $\mathbb{R}^2$ (one component) and $x^2 + y^2 = -1$ (empty, zero components) are the only possibilities. (In $\mathbb{P}^2(\mathbb{R})$, the conic $x^2 + y^2 = z^2$ is one component; $x^2 + y^2 = -z^2$ is empty.)
        \item A smooth real cubic ($d = 3$) has genus 1, so it has at most 2 connected components. The cubic $y^2 = x^3 - x$ has two components: an oval and an unbounded branch. Verify this by plotting the curve. The cubic $y^2 = x^3 + x$ has only one component (no oval). Explain why this is consistent with Harnack's bound.
        \item A smooth real quartic ($d = 4$) has genus 3, so it has at most 4 connected components. Can you find a real quartic with 4 components? (Hint: try $(x^2 + y^2 - 1)(x^2 + y^2 - 4) = 0$---but this is singular. A smooth quartic with 4 ovals is known as a \textbf{M-curve}.)
    \end{enumerate}

\end{enumerate}
